% 表示论 第一章: 代数
% 转录自 scan/reptheory/表示论-ch01.pdf (10 页手稿)

\section{代数}

\subsection{代数}

\begin{definition}[$\mathbb{F}$-代数]\label{def:k-algebra}
给定域 $\mathbb{F}$, 称 $A$ 为 $\mathbb{F}$ 上的代数, 若:
\begin{enumerate}
  \item $A$ 是 $\mathbb{F}$-线性空间;
  \item 存在二元运算 $m\colon A\times A\to A$, $(a,b)\mapsto ab$, 使得
  \begin{enumerate}
    \item 存在 $1\in A$, 使 $1\cdot a=a\cdot 1=a$, $\forall a\in A$;
    \item $\forall a,b,c\in A$, $(a\cdot b)\cdot c=a\cdot(b\cdot c)$;
    \item $m$ 是双线性的, 即
    \begin{align*}
      a\cdot(b+c)&=a\cdot b+a\cdot c, & a\cdot(kb)&=k(a\cdot b),\\
      (a+b)\cdot c&=a\cdot c+b\cdot c, & (ka)\cdot b&=k(a\cdot b),
    \end{align*}
    其中 $a,b,c\in A$, $k\in \mathbb{F}$.
  \end{enumerate}
\end{enumerate}
称 $m$ 为 $A$ 的乘法.
\end{definition}

\begin{example}\label{eg:algebra}
\mbox{}\par\vspace{-0.2em}
\begin{enumerate}
  \item 给定数域 $\mathbb{F}$, 多项式环 $\mathbb{F}[x]$ 是 $\mathbb{F}$-代数.
  \item 给定数域 $\mathbb{F}$, 全体 $n$ 阶方阵 $M_n(\mathbb{F})$ 是 $\mathbb{F}$-代数, 称为全矩阵代数.
  \item 设 $V$ 为 $\mathbb{F}$-线性空间, 则 $\operatorname{End}(V)$ 是 $\mathbb{F}$-代数.
  \item Hamilton 四元数 $\mathbb{H}=\operatorname{span}_{\mathbb{R}}\{1,i,j,k\}$ 是 $\mathbb{R}$-代数, 其乘法表为
  \begin{center}
  \begin{tabular}{c|cccc}
    $\cdot$ & $1$ & $i$ & $j$ & $k$\\
    \hline
    $1$ & $1$ & $i$ & $j$ & $k$\\
    $i$ & $i$ & $-1$ & $k$ & $-j$\\
    $j$ & $j$ & $-k$ & $-1$ & $i$\\
    $k$ & $k$ & $j$ & $-i$ & $-1$
  \end{tabular}
  \end{center}
  \item 群代数. 给定群 $G$, 域 $\mathbb{F}$, 记 $\mathbb{F}G=\operatorname{span}_{\mathbb{F}}\{g \mid g\in G\}$, 其中对 $g,g'\in G$, $g\cdot g'$ 用 $G$ 中的乘法定义为 $gg'\in G$. 可以证明 $\mathbb{F}G$ 为 $\mathbb{F}$-代数.
\end{enumerate}
\end{example}

\begin{example}\label{eg:group-algebra-Z2}
群 $\mathbb{Z}_2=\langle g \mid g^2=1\rangle$. 作同态 $\mathbb{Z}_2\to M_2(\mathbb{F})$:
\[
  1\longmapsto \begin{pmatrix}1&0\\0&1\end{pmatrix},\qquad
  -1\longmapsto \begin{pmatrix}0&1\\1&0\end{pmatrix}.
\]

于是
\[
  \mathbb{F}\mathbb{Z}_2=\operatorname{span}_{\mathbb{F}}\left\{
  \begin{pmatrix}1&0\\0&1\end{pmatrix},
  \begin{pmatrix}0&1\\1&0\end{pmatrix}\right\}
\]

为群代数.
\end{example}

\subsection{子代数}

\begin{definition}[子代数]\label{def:subalgebra}
设 $A$ 为 $\mathbb{F}$-代数, $\varnothing\neq B\subseteq A$, 称 $B$ 为 $A$ 的子代数, 若:
\begin{enumerate}
  \item $1_A\in B$ (单位元必须和 $A$ 中的一样);
  \item $B$ 在 $A$ 的运算下构成 $\mathbb{F}$-代数.
\end{enumerate}
\end{definition}

\begin{remark}\label{rem:unit-in-group}
群的定义中不要求子集的单位元与整个群一致, 因为乘法封闭性已经保证了这一点;
子代数则不同, 必须显式要求 $1_A\in B$.
\end{remark}

\begin{remark}\label{rem:unit-mismatch}
对 Definition \ref{def:subalgebra} 中 (1) 的举例.
$M_n(\mathbb{F})$ 和 $M_{n+1}(\mathbb{F})$ 为 $\mathbb{F}$-代数, 存在嵌入
$\lambda\colon M_n(\mathbb{F})\hookrightarrow M_{n+1}(\mathbb{F})$,
让
\[
  A_{n\times n}\longmapsto \begin{pmatrix}A&0\\0&0\end{pmatrix}.
\]
在嵌入 $\lambda$ 意义下 $M_n(\mathbb{F})\leq M_{n+1}(\mathbb{F})$. 但是 $M_{n+1}(\mathbb{F})$ 的单位元
$I_{n+1}\notin M_n(\mathbb{F})$, 故 $M_n(\mathbb{F})$ 不是 $M_{n+1}(\mathbb{F})$ 的子代数.
\end{remark}

\begin{example}
\mbox{}\par\vspace{-0.2em}
\begin{enumerate}
  \item $\mathbb{F}\subseteq \mathbb{F}[x]$ 为子代数.
  设 $B\subseteq \mathbb{F}[x]$ 为子代数, 且 $1,x\in B$, 那么 $B=\mathbb{F}[x]=\langle x\rangle$
  ($\langle 1,x\rangle=\mathbb{F}[x]$, 但一般默认 $\langle\,\cdot\,\rangle$ 包含 $\mathbb{F}$, 故简记为 $\langle x\rangle$).
  考虑: $\langle x^2\rangle=\operatorname{span}_{\mathbb{F}}\{1,x^2,x^4,\dots,x^{2n},\dots\}$.
\end{enumerate}
\end{example}

\begin{remark}\label{rem:hilbert-basis}
$\mathbb{F}[x]$ 的每个子代数是否都形如 $B=\langle f_1,\dots,f_n\rangle$?
这个问题与 Hilbert 基定理有关, 此处不展开.
\end{remark}

\begin{example}
全矩阵代数 $M_n(\mathbb{F})$ 的子代数:
\begin{enumerate}
  \item 上三角矩阵 (常记为 $B$) 的全体为子代数;
  \item 对角矩阵 (常记为 $T$) 的全体为子代数.
\end{enumerate}
\end{example}

\begin{example}
$\mathbb{F}\mathbb{Z}_2=\operatorname{span}_{\mathbb{F}}\left\{
\begin{pmatrix}1&0\\0&1\end{pmatrix},
\begin{pmatrix}0&1\\1&0\end{pmatrix}\right\}\subset M_2(\mathbb{F})$ 为子代数.
\end{example}

进一步地, 可定义 $\mathbb{Z}_3=\langle g \mid g^3=1\rangle=\{1,g,g^2\}$, 同态 $\mathbb{Z}_3\to M_3(\mathbb{F})$:
\[
  g\longmapsto \begin{pmatrix}0&0&1\\1&0&0\\0&1&0\end{pmatrix},
\]
\[
  \mathbb{F}\mathbb{Z}_3=\operatorname{span}_{\mathbb{F}}\left\{
  \begin{pmatrix}1&0&0\\0&1&0\\0&0&1\end{pmatrix},
  \begin{pmatrix}0&0&1\\1&0&0\\0&1&0\end{pmatrix},
  \begin{pmatrix}0&1&0\\0&0&1\\1&0&0\end{pmatrix}\right\}.
\]

于是, 对一般的 $n$ 令 $\mathbb{Z}_n=\langle g \mid g^n=1\rangle=\{1,g,\dots,g^{n-1}\}$, 并把 $g$ 映到循环置换矩阵
\[
  g\longmapsto
  \begin{pmatrix}
    0&\cdots&0&1\\
    1&\cdots&0&0\\
    \vdots&\ddots&\vdots&\vdots\\
    0&\cdots&1&0
  \end{pmatrix},
\]
即得 $\mathbb{Z}_n$ 的一个矩阵实现.

\begin{example}[子群代数]
设 $H<G$, 则 $\mathbb{F}H\subseteq FG$ 为子代数.
\end{example}

\begin{remark}\label{rem:subalgebra-not-subgroup}
由 $\mathbb{F}H\subseteq FG$ 不能推出 $H<G$. 反例: 取 $\mathbb{F}=\mathbb{C}$, 仍考虑
$\mathbb{Z}_3=\langle g \mid g^3=1\rangle$, $\mathbb{C}\mathbb{Z}_3=\operatorname{span}_{\mathbb{C}}\{1,g,g^2\}$,
记 $\omega=\exp\left(\frac{2\pi i}{3}\right)$.
现考虑元素
\[
  a=\frac{1}{3}\left(1+\omega g+\omega^2 g^2\right)\in\mathbb{C}\mathbb{Z}_3 .
\]

直接计算得
\[
  a^2=\frac{1}{9}\left(1+\omega g+\omega^2 g^2\right)^2
  =\frac{1}{9}\left(3+3\omega g+3\omega^2 g^2\right)
  =\frac{1}{3}\left(1+\omega g+\omega^2 g^2\right)=a .
\]
这说明 $a$ 是幂等的. 取 $B=\operatorname{span}_{\mathbb{C}}\{1,a\}$. 由 $a^2=a$ 知 $B$ 对乘法封闭,
故 $B\subseteq\mathbb{C}\mathbb{Z}_3$ 为子代数, 但 $\{1,a\}$ 不是 $\mathbb{Z}_3$ 的子群.
子代数结构因此比子群丰富.
\end{remark}

\begin{example}
四元数代数 $\mathbb{H}$ 中每个非零元都可逆.
设 $x=a+bi+cj+dk$ ($a,b,c,d$ 不全为 $0$), $\bar{x}=a-bi-cj-dk$, 则
\[
  x^{-1}=\frac{\bar{x}}{|x|^2}=\frac{a-bi-cj-dk}{a^2+b^2+c^2+d^2}.
\]
\end{example}

\begin{definition}[可除代数]\label{def:division-algebra}
称代数 $A$ 是可除代数, 若 $A$ 中每个非零元都有逆元.
\end{definition}

\begin{theorem}[Frobenius]\label{thm:frobenius}
$\mathbb{R}$ 上有限维可除代数只有 $\mathbb{R}$, $\mathbb{C}$, $\mathbb{H}$.
\end{theorem}

\subsection{子代数的交}

\begin{proposition}\label{prop:intersection}
设 $B_1,B_2\subseteq A$ 为 $A$ 的子代数, 则 $B_1\cap B_2$ 仍为 $A$ 的子代数.
\end{proposition}

\begin{proof}
$1_A$ 同时属于 $B_1$ 与 $B_2$, 故 $1_A\in B_1\cap B_2$. 设 $x,y\in B_1\cap B_2$, $k\in F$.
因为 $B_1,B_2$ 都是子代数, 所以 $x+y$, $kx$, $xy$ 同时属于 $B_1$ 与 $B_2$, 即属于 $B_1\cap B_2$.
于是 $B_1\cap B_2$ 在 $A$ 的运算下构成 $\mathbb{F}$-代数, 且与 $A$ 有相同的单位元, 故为 $A$ 的子代数.
\end{proof}

\subsection{代数直和}

\begin{definition}[外直和]\label{def:direct-sum}
设 $A_1,A_2$ 为 $\mathbb{F}$-代数, $A_1,A_2$ 的外直和记作 $A_1\times A_2$, 定义为 $A_1,A_2$ 的 Cartesian 积:
\[
  A_1\times A_2=\{(a_1,a_2)\mid a_i\in A_i,\ i=1,2\},
\]
直和上运算定义如下:
\begin{align*}
  \text{加法: } & (a_1,a_2)+(a_1',a_2')=(a_1+a_1',a_2+a_2'),\\
  \text{数乘: } & k\cdot(a_1,a_2)=(ka_1,ka_2),\\
  \text{乘法: } & (a_1,a_2)\cdot(a_1',a_2')=(a_1a_1',a_2a_2'),
\end{align*}
其中 $k\in F$, $a_i,a_i'\in A_i$. 直和的单位元为 $(1_{A_1},1_{A_2})$.
\end{definition}

通常可将 $(a_1,a_2)$ 写作 $a_1+a_2$ (不是代数上的加法). 这种写法蕴含内直和的思想,
将 $A_1,A_2$ 视为 $A_1\times A_2$ 的子空间, 即可构造嵌入
$\lambda\colon A_1\hookrightarrow A_1\times A_2$, $a_1\mapsto(a_1,0)$, 对 $A_2$ 类似.
内直和 (记作 $A_1\oplus A_2$) 需满足:
\begin{enumerate}
  \item $\forall a\in A_1\oplus A_2$, $\exists!\, a_1\in A_1$, $a_2\in A_2$, $a=a_1+a_2$;
  \item $A_1\cap A_2=\{0\}$.
\end{enumerate}
在这种记法下, 有
\[
  a_1+a_2=(a_1,0)+(0,a_2)=(a_1,a_2),\qquad
  a_1\cdot a_2=(a_1,0)\cdot(0,a_2)=(0,0)=0 .
\]

内直和在时, 内外直和没有本质区别, 故统一记作 $A_1\oplus A_2$, 而不加区分.

\begin{remark}\label{rem:direct-sum-not-subalgebra}
$A_1,A_2$ 都不是 $A_1\oplus A_2$ 的子代数. 按 Definition \ref{def:subalgebra} 的第 (1) 条, 子代数必须与 $A_1\oplus A_2$ 有相同的单位元; 而 $A_1$ 的单位元是 $(1,0)$, $A_1\oplus A_2$ 的单位元是 $(1,1)$, 二者不同.
\end{remark}

\begin{example}
\mbox{}\par\vspace{-0.2em}
\begin{enumerate}
  \item 设 $M_{n_1}(\mathbb{F})$, $M_{n_2}(\mathbb{F})$ 均为 $\mathbb{F}$-代数, 则
  \[
    M_{n_1}(\mathbb{F})\oplus M_{n_2}(\mathbb{F})
    =\left\{\begin{pmatrix}A&0\\0&B\end{pmatrix}
    \ \middle|\ A\in M_{n_1}(\mathbb{F}),\ B\in M_{n_2}(\mathbb{F})\right\}.
  \]
  \item 代数 $\mathbb{F}[x]$. 设 $B_1=\mathbb{F}[x^2]$, $B_2=\mathbb{F}[x^3]$, 则 $B_1,B_2$ 均为 $\mathbb{F}[x]$ 的子代数,
  空间的和 $B_1+B_2$ 构成代数 $\mathbb{F}\langle x^2,x^3\rangle$.
  $B_1,B_2$ 没有内直和, 因为 $B_1\cap B_2=\mathbb{F}[x^6]\neq\{0\}$.
  此外, 外直和 $B_1\times B_2$ 与 $\mathbb{F}\langle x^2,x^3\rangle$ 不同构,
  因为对 $b_1\in B_1$, $b_2\in B_2$ 有
  $(b_1,0)\cdot(0,b_2)=0$,
  这意味着 $B_1\times B_2$ 有零因子.
\end{enumerate}
\end{example}

\subsection{代数理想}

\begin{definition}[理想]\label{def:ideal}
设 $A$ 为 $\mathbb{F}$-代数, 称子空间 $I\subseteq A$ 是 $A$ 的左 (右) 理想, 若 $AI\subseteq I$ ($IA\subseteq I$).
若 $I$ 既是左理想, 又是右理想, 则称 $I$ 是理想.
左理想记作 $I\trianglelefteq_l A$, 右理想记作 $I\trianglelefteq_r A$, 理想记作 $I\trianglelefteq A$.
\end{definition}

左右的区分: 理想是 "吸收" $A$ 中元素, 从左吸收即在乘 $A$ 中元素, 类似定义 "右".

\begin{remark}\label{rem:summand-ideal}
$A_1,A_2$ 都是 $A_1\oplus A_2$ 的理想.
\end{remark}

\begin{proof}
只证 $A_1$. 设 $a_1,a_1'\in A_1$, $a_2\in A_2$, 有
\begin{align*}
  (a_1+a_2)\cdot a_1' &= a_1\cdot a_1'+a_2\cdot a_1'=a_1\cdot a_1'+0=a_1a_1'\in A_1,\\
  a_1'\cdot(a_1+a_2) &= a_1'\cdot a_1+a_1'\cdot a_2=a_1'a_1+0=a_1'a_1\in A_1 .
\end{align*}
从而 $(A_1\oplus A_2)\cdot A_1\subseteq A_1$, 且 $A_1\cdot(A_1\oplus A_2)\subseteq A_1$.
故 $A_1\trianglelefteq(A_1\oplus A_2)$.
\end{proof}

\begin{example}
\mbox{}\par\vspace{-0.2em}
\begin{enumerate}
  \item 考虑全矩阵代数 $M_n(\mathbb{F})$. 记 $E_{ij}$ 为第 $(i,j)$ 个矩阵单位, $e_j$ 为第 $j$ 个标准基向量.
  设 $I_l\trianglelefteq_l M_n(\mathbb{F})$ 为左理想, 令
  \[
    W=\{Ae_j \mid A\in I_l,\ 1\le j\le n\}\subseteq \mathbb{F}^{n}
  \]
  为 $I_l$ 中所有矩阵的列向量张成的子空间, 则
  \[
    I_l=\{A\in M_n(\mathbb{F}) \mid \operatorname{Im}(A)\subseteq W\}.
  \]
  事实上, "$\subseteq$" 由 $W$ 的定义直接得到. 反之设 $A$ 的每一列都属于 $W$,
  则对每个 $j$ 存在 $B\in I_l$ 与 $k$ 使 $Ae_j=Be_k$, 于是
  $Ae_j e_j^{\mathrm{T}}=B e_k e_j^{\mathrm{T}}=BE_{kj}\in I_l$;
  由 $A=\sum_j Ae_j e_j^{\mathrm{T}}$ 得 $A\in I_l$.

  于是左理想与 $\mathbb{F}^{n}$ 的子空间一一对应. 取 $W=\operatorname{span}_{\mathbb{F}}\{e_{i_1},\dots,e_{i_k}\}$,
  即得"只有第 $i_1,\dots,i_k$ 列可能非零"的左理想, 其元素形如 (其中 $0$ 为零列)
  \[
    \left(\alpha_{i_1}\ \ 0\ \alpha_{i_2}\ \ 0\ \cdots\ 0\ \alpha_{i_k}\ \ 0\ \cdots\right),
    \qquad \alpha_{i_j}\in \mathbb{F}^{n\times 1};
  \]
  取 $W=\mathbb{F}^{n}$ 得 $I_l=M_n(\mathbb{F})$, 取 $W=\{0\}$ 得 $I_l=\{0\}$.

  同理, 对右理想 $I_r\trianglelefteq_r M_n(\mathbb{F})$, 令 $U$ 为其所有行向量张成的子空间, 则
  $I_r=\{A\in M_n(\mathbb{F})\mid$ $A$ 的每一行都属于 $U\}$;
  取 $U=\operatorname{span}_{\mathbb{F}}\{e_{j_1}^{\mathrm{T}},\dots,e_{j_k}^{\mathrm{T}}\}$ 即得"只有第 $j_1,\dots,j_k$ 行可能非零"的右理想,
  其元素形如 (其中 $0$ 为零行)
  \[
    \begin{pmatrix}
      \beta_{j_1}\\ 0\\ \beta_{j_2}\\ 0\\ \vdots\\ \beta_{j_k}\\ 0
    \end{pmatrix},
    \qquad \beta_{j_i}\in \mathbb{F}^{1\times n}.
  \]

  最后设 $I\trianglelefteq M_n(\mathbb{F})$ 同时为左、右理想. 若 $I\neq\{0\}$, 取 $0\neq A\in I$
  并设 $a_{ij}\neq 0$, 则 $E_{ij}=a_{ij}^{-1}E_{ii}AE_{jj}\in I$;
  于是全部矩阵单位都属于 $I$, 从而 $I=M_n(\mathbb{F})$.
  故 $I$ 只能是 $\{0\}$ 或 $M_n(\mathbb{F})$.
\end{enumerate}
\end{example}

\begin{definition}[单代数]\label{def:simple-algebra}
称代数 $A$ 是单代数, 若 $A$ 没有非平凡理想 (即只有 $\{0\}$ 和 $A$ 本身).
\end{definition}

\begin{theorem}\label{thm:simple-over-C}
$\mathbb{C}$ 上所有有限维单代数只有 $M_n(\mathbb{C})$.
\end{theorem}

\begin{example}
$\mathbb{H}$ 是单代数.
设 $I\trianglelefteq\mathbb{H}$, 若 $I\neq\{0\}$, 取 $0\neq x\in I$.
既然 $\mathbb{H}$ 是可除代数, 故 $\exists x^{-1}\in\mathbb{H}$, 有 $x^{-1}x=x\cdot x^{-1}=1\in I$,
$\Rightarrow I=\mathbb{H}$. 于是 $\mathbb{H}$ 的理想只有 $\{0\}$ 和 $\mathbb{H}$, 故为单代数.
\end{example}

\begin{example}
回忆 $\mathbb{Z}_3=\langle g \mid g^3=1\rangle=\{1,g,g^2\}$,
$\mathbb{F}\mathbb{Z}_3=\operatorname{span}_{\mathbb{F}}\{1,g,g^2\}$.
令 $\mathbb{F}=\mathbb{C}$, 仍取 $a=\frac{1}{3}(1+\omega g+\omega^2 g^2)$, 其中 $\omega=\exp\left(\frac{2\pi i}{3}\right)$,
则 $I=\operatorname{span}_{\mathbb{F}}\{a\}=\operatorname{span}_{\mathbb{C}}\{a\}\subseteq\mathbb{C}\mathbb{Z}_3$ 为子空间 (不是子代数!).
下面说明 $I\trianglelefteq\mathbb{C}\mathbb{Z}_3$. 由于 $\mathbb{Z}_3$ 交换, 只需证
$\mathbb{C}\mathbb{Z}_3\cdot I\subseteq I$; 又只需对基验证.
\begin{align*}
  g^{\phantom{2}}\cdot a &= g\cdot\frac{1}{3}\left(1+\omega g+\omega^2 g^2\right)
  =\frac{1}{3}\left(g+\omega g^2+\omega^2\right)
  =\frac{1}{3}\cdot\omega^2\left(1+\omega g+\omega^2 g^2\right)=\omega^2 a\in I,\\
  g^2\cdot a &= g\cdot(g\cdot a)=\omega^2\cdot g\cdot a=\omega^4 a=\omega a\in I .
\end{align*}
故 $I\trianglelefteq\mathbb{C}\mathbb{Z}_3$.
\end{example}

\subsection{商代数}

给定代数 $A$, 子空间 $I$. 自然 $A$ 是线性空间, 那么便可以构成商空间 $A/I$,
其中 $A/I=\{[a]\mid a\in A\}$, $[a]$ 为 $a$ 的等价类, 即 $\forall a,b\in A$,
$[a]=[b]\Leftrightarrow a-b\in I$. 加法与数乘在商空间 $A/I$ 上自然良定义.
若希望 $A/I$ 为代数, 则要在其上定义乘法. 一个自然的想法是定义
\[
  A/I\times A/I\longrightarrow A/I\colon\quad ([a],[b])\longmapsto[ab].
\]

需要验证该定义合理, 即运算结果不依赖代表元的选取. 这就要求 $I$ 为理想.
我们说明: $I$ 为理想时, 该运算是 well-defined 的.

\begin{proof}
只证一边. 设 $[a]=[a']$, 要证 $[ab]=[a'b]$.
由假设有 $a-a'\in I$, 右理想条件有 $(a-a')b\in I$ ($b\in A$), 则
$ab-a'b\in I\Rightarrow[ab]=[a'b]$.
\end{proof}

\begin{definition}[商代数]\label{def:quotient-algebra}
称在上述乘法定义下形成的代数 $A/I$ 为 $A$ 模去 $I$ 形成的商代数.
\end{definition}

\begin{remark}\label{rem:ideal-necessary}
$I$ 为理想是必要的, 我们用一个例子来说明:
取全矩阵代数 $M_2(\mathbb{F})$, 设
$I=\left\{\begin{pmatrix}0&x\\0&0\end{pmatrix}\ \middle|\ x\in F\right\}$.
显然, $I$ 为子空间而不是理想.
\end{remark}

设
\[
  a=\begin{pmatrix}0&0\\1&0\end{pmatrix},\qquad
  i=\begin{pmatrix}0&1\\0&0\end{pmatrix}
\]
在商空间 $M_2(\mathbb{F})/I$ 中.
由于 $i\in I$, 故 $[i]=[0]$. 按定义计算两个乘法:
\begin{align*}
  [a]\cdot[i] &= [a\cdot i]
  =\left[\begin{pmatrix}0&0\\1&0\end{pmatrix}\begin{pmatrix}0&1\\0&0\end{pmatrix}\right]
  =\left[\begin{pmatrix}0&0\\0&1\end{pmatrix}\right],\\
  [a]\cdot[0] &= [a\cdot 0]=[0].
\end{align*}
显然 $\left[\begin{pmatrix}0&0\\0&1\end{pmatrix}\right]\neq[0]$, 则乘法结果依赖于代表元的选取,
此时乘法不是 well-defined 的.

\begin{example}
设 $A_1,A_2$ 为 $\mathbb{F}$-代数, 记 $A=A_1\oplus A_2$, 那么 $A/A_1\cong A_2$, $A/A_2\cong A_1$.
\end{example}

\begin{proof}
只证明 $A/A_1\cong A_2$. 为方便, 写对 $\forall a=a_1+a_2$, $a_{1,2}\in A_{1,2}$,
记 $\pi_2a=a_2$. 取 $[a],[b]\in A/A_1$, 若 $[a]=[b]$, 那么 $a-b\in A_1$.
注意此时的 $A_1$ 是嵌入意义下的, 即 $A_1$ 中元素形如 $(a_1,0)$, 于是 $\pi_2a=\pi_2b$.
写 $\varphi\colon A/A_1\to A_2$, $[a]\mapsto\pi_2a$.
由上面的说明 $\varphi$ 是单射; 又 $\varphi$ 显然为满射, 故 $\varphi$ 为双射.
又:
\[
  \varphi([a][b])=\varphi([ab])=\pi_2(ab)=\pi_2(a)\pi_2(b)=\varphi([a])\varphi([b]),
  \quad\forall a,b\in A .
\]
故 $\varphi$ 是同态, 从而 $\varphi$ 为同构, 有 $A/A_1\cong A_2$.
\end{proof}

\subsection{代数同态}

\begin{definition}[代数同态]\label{def:algebra-hom}
设 $A,B$ 为 $\mathbb{F}$-代数, 称线性映射 $f\colon A\to B$ 为同态, 若:
\begin{enumerate}
  \item $\forall a,b\in A$, $f(ab)=f(a)f(b)$;
  \item $f(1_A)=1_B$.
\end{enumerate}
若 $f$ 是单射, 则称 $f$ 是单同态. 若 $f$ 是满射, 则称 $f$ 是满同态.
若 $f$ 是双射, 则称 $f$ 是同构, 此时记 $A\cong B$.
\end{definition}

\begin{example}
\mbox{}\par\vspace{-0.2em}
\begin{enumerate}
  \item 线性空间 $V$, $\dim(V)=n$, $V$ 有一组基 $\{\alpha_1,\dots,\alpha_n\}$.
  $V$ 上线性变换 $\sigma\in\operatorname{End}(V)$, 给出
  $\sigma(\alpha_1,\dots,\alpha_n)=(\alpha_1,\dots,\alpha_n)A$, $A\in M_n(\mathbb{F})$.
  $\sigma\mapsto A_\sigma$ 是代数同态 (实际上 $\operatorname{End}(V)\cong M_n(\mathbb{F})$).
  \item $f\colon \mathbb{F}\mathbb{Z}_2\to M_2(\mathbb{F})$,
  \[
    1\mapsto\begin{pmatrix}1&0\\0&1\end{pmatrix},\qquad
    g\mapsto\begin{pmatrix}0&1\\1&0\end{pmatrix}
  \]
  为代数同态.
\end{enumerate}
\end{example}

\begin{definition}[典范同态]\label{def:canonical-hom}
设 $A$ 为代数, $I\trianglelefteq A$, 映射 $\pi\colon A\to A/I$, $a\mapsto[a]$, 称 $\pi$ 为典范同态.
\end{definition}

\begin{example}
设 $A=\mathbb{F}[x]$, 任取 $a_0\in \mathbb{F}$, 记理想 $I_{a_0}=(x-a_0)$.
则 $\pi_{a_0}\colon \mathbb{F}[x]\to \mathbb{F}[x]/I_{a_0}$ 为典范同态.
可以考察元素的像, 以 $x$, $x^2$ 为例.
$x\mapsto[x]$, 由于 $I_{a_0}=(x-a_0)$, 则 $x-a_0\in I_{a_0}\Rightarrow[x]=[a_0]$.
$x^2\mapsto[x^2]$, 则 $[x^2]=[x]^2=[a_0]^2=[a_0^2]$.
\end{example}

\begin{theorem}[同态基本定理]\label{thm:first-iso}
设 $f\colon A\to B$ 为代数同态, 则
\begin{enumerate}
  \item $\operatorname{Ker}(f)=\{a\in A\mid f(a)=0\}\trianglelefteq A$;
  \item $\operatorname{Im}(f)=\{b\in B\mid\exists a\in A,\ \text{s.t.}\ f(a)=b\}\subseteq B$
  为 $B$ 的子代数;
  \item $A/\operatorname{Ker}(f)\cong\operatorname{Im}(f)$.
\end{enumerate}
\end{theorem}

\begin{proof}
\begin{enumerate}
\item $\forall a\in A$, $a_0\in\operatorname{Ker}(f)$, 有
\begin{align*}
  f(a\cdot a_0) &= f(a)\cdot f(a_0)=f(a)\cdot 0=0
  \Rightarrow a\cdot a_0\in\operatorname{Ker}(f)
  \Rightarrow A\cdot\operatorname{Ker}(f)\subseteq\operatorname{Ker}(f),\\
  f(a_0\cdot a) &= f(a_0)\cdot f(a)=0\cdot f(a)=0
  \Rightarrow a_0\cdot a\in\operatorname{Ker}(f)
  \Rightarrow \operatorname{Ker}(f)\cdot A\subseteq\operatorname{Ker}(f).
\end{align*}

于是 $\operatorname{Ker}(f)\trianglelefteq A$.

\item $\operatorname{Im}(f)\subseteq B$ 是显然的, 只要说明 $\operatorname{Im}(f)$ 是代数.
即 $\operatorname{Im}(f)$ 对加法, 数乘, 乘法封闭;
其余性质由 $B$ 中自然继承.
取 $b_{1,2}\in\operatorname{Im}(f)$, 则 $\exists a_{1,2}\in A$, s.t. $f(a_{1,2})=b_{1,2}$.

由于 $f$ 是线性映射 (参看 Definition \ref{def:algebra-hom}), 故
\begin{align*}
  b_1+b_2 &= f(a_1)+f(a_2)=f(a_1+a_2)\in\operatorname{Im}(f),\\
  kb_1 &= kf(a_1)=f(ka_1)\in\operatorname{Im}(f),
\end{align*}
其中 $k\in \mathbb{F}$. 其次 $b_1b_2=f(a_1)f(a_2)=f(a_1a_2)\in\operatorname{Im}(f)$.
故 $\operatorname{Im}(f)$ 是 $B$ 的子代数.

\item 定义映射 $\varphi\colon A/\operatorname{Ker}(f)\to\operatorname{Im}(f)$, $[a]\mapsto f(a)$.
下面说明 $\varphi$ 是同构. 先说明 $\varphi$ 是同态:
$\varphi([1])=f(1_A)=1_B$, 且
\[
  \varphi([a][b])=\varphi([ab])=f(ab)=f(a)f(b)=\varphi([a])\varphi([b]),
\]
其中 $a,b\in A$. 故 $\varphi$ 是同态. 下证 $\varphi$ 是双射.

\begin{enumerate}
\item $\varphi$ 是单射. 设 $f(a_1)=f(a_2)\Rightarrow f(a_1-a_2)=0
\Rightarrow a_1-a_2\in\operatorname{Ker}(f)\Rightarrow[a_1]=[a_2]$. 得证.

\item $\varphi$ 是满射. $\forall b\in\operatorname{Im}(f)$, $\exists a\in A$, $f(a)=b$.
故总有 $[a]\in A/\operatorname{Ker}(f)$, s.t. $\varphi([a])=f(a)=b$.

\end{enumerate}
\end{enumerate}

综合以上, $\varphi$ 是同构, 故 $A/\operatorname{Ker}(f)\cong\operatorname{Im}(f)$.
\end{proof}

\begin{example}
\mbox{}\par\vspace{-0.2em}
\begin{enumerate}
  \item 设 $A=\mathbb{F}[x]$. $\forall a\in \mathbb{F}$, 记同态 $f_a\colon \mathbb{F}[x]\to \mathbb{F}[a]$, $p(x)\mapsto p(a)$.
  于是 $\operatorname{Im}(f_a)=\mathbb{F}[a]$, $\operatorname{Ker}(f_a)=(x-a)$, 从而
  \[
    \mathbb{F}[x]/(x-a)\cong \mathbb{F}[a].
  \]
\end{enumerate}
\end{example}
